\chapter{Gum Swelling}

\section*{Definition}
Localized enlargement of gum tissue resulting from inflammation, edema, accumulation of fluid and cellular components, vascular changes at the gingival margin area surrounding the tooth neck; visible protrusion beyond normal contour, typically red, tender, bleeding upon brushing, eating, or spicy foods.

\section*{Pathophysiology}
Gum swelling arises from inflammatory response triggered by bacterial plaque, biofilm, dental caries, calculus/tartar, foreign material, trauma, ill-fitting restorations, braces, retainers, dentures, implants, systemic medications, immunosuppressants, hormonal fluctuations, pregnancy, puberty, thyroid dysfunction, blood disorders, hematologic diseases, leukemias, lymphomas, autoimmune conditions, lupus, vasculitides, granulomatous diseases, Crohn's, sarcoidosis, Behçet's, drug eruptions, angioedema, allergic reactions, infections (viral: herpes simplex HSV, Epstein-Barr EBV, cytomegalovirus CMV, HIV), HIV-associated periodontitis, necrotizing ulcerative gum disease (NUG/ANUG), acute necrotizing ulcerative gingivitis, Vincent infection (spirochetes, Fusobacterium necrophorum), anaerobic bacteria overgrowth.

\section*{Etiology}
\begin{itemize}
  \item Gingivitis (early reversible stage of gum disease)
  \item Periodontal disease (gum disease progressing to bone loss)
  \item Dental abscess (infection at tooth root or between gums)
  \item Impacted wisdom teeth
  \item Trauma or injury to gums
  \item Poor oral hygiene
  \item Hormonal changes (pregnancy, puberty, menstruation)
  \item Certain medications (phenytoin, cyclosporine, calcium channel blockers)
  \item Vitamin deficiencies (especially vitamin C)
  \item Blood disorders (leukemia, anemia)
\end{itemize}

\section*{Indications for Evaluation}
Seek care if:
\begin{itemize}
  \item Severe or worsening swelling
  \item Difficulty breathing or swallowing
  \item Fever with gum swelling
  \item Persistent bleeding
  \item Swelling interfering with speech or eating
  \item Loose teeth without obvious cause
\end{itemize}

\section*{Related Symptoms}
\begin{itemize}
  \item Bad breath (halitosis)
  \item Tooth pain
  \item Mouth odor
  \item Difficulty chewing
  \item Gum tenderness
  \item Swollen lymph nodes in neck
\end{itemize}
\textbf{Note:} Always consider serious underlying conditions when evaluating persistent gum swelling.

\section*{Self-Care / Home Management}
\begin{itemize}
  \item Maintain good oral hygiene: brush twice daily, floss daily, and rinse with warm salt water (1/2 tsp salt in 8 oz water).
  \item Avoid tobacco, alcohol, and spicy or acidic foods that may irritate oral tissues.
  \item Stay hydrated by sipping water throughout the day to combat dry mouth.
  \item Over-the-counter pain relievers and topical oral gels can provide temporary relief for toothache or gum discomfort.
  \item See a dentist promptly for persistent oral symptoms, as early intervention improves outcomes.
\end{itemize}

\section*{Diagnostic Approach}
\begin{itemize}
  \item History covers onset, character of pain, aggravating/relieving factors, oral hygiene habits, and dental care history.
  \item Oral examination inspects teeth, gums, mucosa, tongue, and oropharynx for caries, inflammation, lesions, or masses.
  \item Dental X-rays (panoramic, periapical) identify caries, abscesses, impactions, and bone pathology.
  \item Referral to a dentist or oral surgeon is appropriate for definitive diagnosis and treatment of dental pathology.
\end{itemize}

\section*{Treatment / Management Overview}
\begin{itemize}
  \item Dental caries require restorative treatment (fillings, crowns, root canal) by a dentist.
  \item Oral infections (abscess, pericoronitis) may require drainage, antibiotics, and definitive dental treatment.
  \item Dry mouth management includes hydration, sugar-free lozenges, saliva substitutes, and addressing underlying causes.
  \item Regular dental check-ups every 6 months are essential for prevention and early detection of oral pathology.
\end{itemize}

\clearpage
